\documentclass[14pt]{extarticle}
\usepackage{preamble}
\geometry{letterpaper, landscape, margin=0.5in}
\usepackage{qrcode}
\renewcommand{\answer}[1]{}
\setlength{\parskip}{4pt}

\begin{document}
\pagestyle{empty}

\daybanner{Unit 2 \textperiodcentered\ Topic 2.3 \textperiodcentered\ TODAY'S PROBLEM}{Can Chance Make a Streak?}

\noindent\begin{minipage}[t]{0.57\linewidth}
\vspace{0pt}
Suppose a coin advertised as fair is flipped 10 times and gives HHHHHHTTTT. To investigate what chance alone could produce, a class simulates 100 sets of 10 independent fair-coin flips. For each set, the class records the longest run of matching outcomes; the dotplot shows the results.\par\smallskip \textbf{Claim A:} ``The coin is rigged---six heads in a row cannot be chance.''\par \textbf{Claim B:} ``Streaks like this can happen by chance; simulation can tell us how unusual this one is.''\par
\end{minipage}\hfill
\begin{minipage}[t]{0.40\linewidth}
\vspace{0pt}
\begin{center}
\begin{tikzpicture}
\begin{axis}[
  width=0.85\linewidth,
  height=1.60in,
  ymin=0, ymax=33, ytick=\empty, axis y line=none, axis x line=bottom,
  xlabel={Longest run in one set of 10 fair-coin flips}, tick label style={font=\small}, label style={font=\small}
]
\addplot[only marks, mark=*, mark size=2.4pt, color=dayblue] coordinates {
  (2,1)
  (2,2)
  (2,3)
  (2,4)
  (2,5)
  (2,6)
  (2,7)
  (2,8)
  (2,9)
  (2,10)
  (2,11)
  (2,12)
  (2,13)
  (2,14)
  (2,15)
  (2,16)
  (2,17)
  (2,18)
  (2,19)
  (2,20)
  (2,21)
  (2,22)
  (3,1)
  (3,2)
  (3,3)
  (3,4)
  (3,5)
  (3,6)
  (3,7)
  (3,8)
  (3,9)
  (3,10)
  (3,11)
  (3,12)
  (3,13)
  (3,14)
  (3,15)
  (3,16)
  (3,17)
  (3,18)
  (3,19)
  (3,20)
  (3,21)
  (3,22)
  (3,23)
  (3,24)
  (3,25)
  (3,26)
  (3,27)
  (3,28)
  (3,29)
  (3,30)
  (3,31)
  (3,32)
  (4,1)
  (4,2)
  (4,3)
  (4,4)
  (4,5)
  (4,6)
  (4,7)
  (4,8)
  (4,9)
  (4,10)
  (4,11)
  (4,12)
  (4,13)
  (4,14)
  (4,15)
  (4,16)
  (4,17)
  (4,18)
  (4,19)
  (4,20)
  (4,21)
  (4,22)
  (4,23)
  (4,24)
  (4,25)
  (5,1)
  (5,2)
  (5,3)
  (5,4)
  (5,5)
  (5,6)
  (5,7)
  (5,8)
  (5,9)
  (5,10)
  (5,11)
  (5,12)
  (5,13)
  (5,14)
  (6,1)
  (6,2)
  (6,3)
  (6,4)
  (6,5)
  (6,6)
  (6,7)
};
\end{axis}
\end{tikzpicture}
\end{center}
\end{minipage}

\begin{firsttakebox}{FIRST TAKE (5 min, on your sheet)}
Which claim seems better supported right now? Cite one detail from the sequence or dotplot.
\end{firsttakebox}

\begin{description}[leftmargin=1.15in, labelwidth=1in, itemsep=2pt, topsep=2pt]
  \item[\dokbadge{1}~(a)] State the longest run in the observed sequence and the most common longest-run value among the simulated sets.
  \item[\dokbadge{2}~(b)] Use the dotplot to estimate the probability that 10 fair-coin flips have a longest run at least as long as the observed run. Interpret the estimate in context.
  \item[\dokbadge{3}~(c)] Which claim is better supported by this evidence? Cite the simulation result that settles your choice. Then describe a larger simulation that would more precisely settle the disagreement: state what one trial is, what to record and count, how many trials to use, and what result would make you doubt that the coin is fair.
\end{description}

\vfill
\noindent\begin{minipage}{0.84\linewidth}
\textbf{Video + follow-along:} \href{https://robjohncolson.github.io/apstats-live-worksheet/u4_lesson1-2_live.html}{\texttt{\detokenize{u4_lesson1-2_live.html}}} (scan the code)\par
\textbf{Finish (a)--(c) after the video. Turn in your sheet.}
\end{minipage}\hfill
\qrcode[height=0.75in]{https://robjohncolson.github.io/apstats-live-worksheet/u4_lesson1-2_live.html}

\end{document}
